Let be $T_{\mathrm{sol}}$ the duration of the solar day, $T_{\mathrm{sid}}$
the duration of the sidereal day and $T_{\mathrm{year}}$ the length of the
year. In the present situation, we have
\[
  \frac{1}{T_{\mathrm{sol}}} = \frac{1}{T_{\mathrm{sid}}}
    - \frac{1}{T_{\mathrm{year}}}
\]
\hfill(3 pt calc.\ $T_{\mathrm{sid}}$)

Using the given data,
$T_{\mathrm{sid}} = 0.997270\ \mathrm{days}
= 23\,\mathrm{h}\ 56\,\mathrm{m}\ 4\,\mathrm{s}$
\hfill(3 pt sidereal doesn't change)

The sidereal day depends only of the magnitude of Earth's rotation speed, and
the angular velocity of the sun's annual movement is opposite. So we have
\[
  \frac{1}{T'_{\mathrm{sol}}} = \frac{1}{T_{\mathrm{sid}}}
    + \frac{1}{T_{\mathrm{year}}}
\]
\hfill(2 pt $-$ to $+$)
\[
  \therefore\ T'_{\mathrm{sol}} = 0.994555\ \mathrm{d}
    = 23\,\mathrm{h}\ 52\,\mathrm{m}\ 9.5\,\mathrm{s}
\]
\hfill(2 pt right number)
